\section{Conclusion}

We presented OMU-MAE, a voxel-level masked autoencoder that uses frozen
DINOv2 features as the cross-modal reconstruction target. Under a five-way
controlled ablation on SemanticKITTI linear probing, OMU-MAE reaches
$19.68\%$ mIoU at $100\%$ labels and wins at every label fraction:
$+14.19$\,pp over random init, $+1.77$\,pp over a faithful Occupancy-MAE
re-implementation (the cross-modal target), $+3.27$\,pp over a no-mask
CleverDistiller-equivalent (the masking bias), and $+4.22$\,pp over a
re-implemented SLidR contrastive baseline (masked reconstruction vs.\
contrastive distillation of the same frozen-VFM target). A nuScenes
cross-sensor transfer study shows the representation carries some
sensor-agnostic signal---it beats random init, Occupancy-MAE and SLidR
out of domain---but the in-domain advantage of masking does not transfer,
a nuance we report honestly. Together with a preliminary negative result on
semantics-driven masking, these experiments isolate \emph{when} each
ingredient of cross-modal masked voxel pretraining helps, and when it does
not. The full pipeline trains in a few hours on a single GPU, and we
release code and pretrained weights to support reproduction.
